\scanpage{107}

\[
\left\{\begin{aligned}
 \mathrm{T3}\ \text{a)}\quad &\forall f\in F\text{ et }s\in S\smallsetminus f,\ \text{il existe une \deleted{suite}\inserted{``chaîne''} de}\\
 &\text{faces }f_0=f,\ f_1,\ldots,f_N\ \text{(dont deux consécutives}\\
 &\text{sont adjacentes) avec }s\in f_N\\
 \text{b)}\quad &\text{Si }\deleted{\{s,t\},\{s,t'\}}\inserted{f,f'\in F}\text{ sont deux \deleted{arêtes adjacentes,}\inserted{faces contenant}}\\
 &\text{un \unsure{$\simeq$|$=$} sommet }s,\\
 &\text{qui peuvent être jointes par une ``chaîne de}\\
 &\text{faces''\ }f_0\inserted{=f},\ldots,f_N\inserted{=f'},\ \text{alors on peut $\simeq$ trouver}\\
 &\text{une telle chaîne formée de }f_i\text{ contenant }s.
\end{aligned}\right.\tag*{(18)}
\]

Une structure $(S,F)$ satisfaisant T1 à T3 peut aussi s'interpréter comme une structure

\[
\begin{aligned}
 \bigl(S,\underbrace{(F_i)_{i\in\Pi}}_{\substack{\text{notation}\\\text{indicielle}\\\text{pour un ensemble }\Pi\\\text{d'objets }F_i}}\bigr)\qquad
 &\text{où }\Pi\subset\unsure{\mathscr P|\mathcal P|\mathfrak P}(\mathscr P_3(S))\\
 &\text{est un ensemble}\\
 &\text{formé de parties}\\
 &F_i\subset\mathscr P_3(S)
\end{aligned}\tag*{(19)}
\]

avec le seul axiome que chaque $F_i$ définit sur $S$ une structure de 2-polyèdre triangulé \inserted{\unsure{o-|0-}}connexe [tout $a\in\mathscr P_2(S)$ contenu dans un $f\in F_i$ est contenu dans exactement deux tels $f$, plus axiome de \unsure{o-connexité|0-connexité} dans tout \deleted{point}\inserted{sommet} $s\in S$, plus axiome de connexité globale\ldots] --

% slanted
\margin{\uline{NB} $\Pi\ne\varnothing\iff S\ne\varnothing$ Pour la suite, on suppose $S\ne\varnothing$, ce qui implique que $\operatorname{card}S\geq4$ (le cas $\operatorname{card}S=4$ est intéressant, c'est celui du tétraèdre (ici $\operatorname{card}\Pi=1$, $f=\mathscr P_3(S)\ldots$).}

Les axiomes qui suivent sont destinés, entre autres, à assurer que \deleted{$\unsure{\Pi\subset\mathscr P}$} l'ens $\Pi$ de structures polyédrales sur $S$ est déterminé canoniquement par la connaissance de l'\uline{une} de ces
% archive sheet number: 106
